\documentclass[12pt]{article}
\usepackage[utf8]{inputenc}
\usepackage[T1]{fontenc}
\usepackage{hyperref}
\title{Usan Language Resources for K-12 Curriculum Development}
\author{Compiled by Claude Sonnet 5.5 (claude-sonnet-5-5) for Andr\'as Kornai}
\date{}
\begin{document}
\maketitle

\section{Introduction}
Usan (ISO 639-3 wnu; also called Wanuma) is a Papuan language of Madang Province, Papua New Guinea, spoken along the southern slopes of the Adelbert Range. Wikipedia gives about 1,400 speakers (a 1991 figure), so the number is probably dated. Glottolog places it in the Madang group, under Croisilles and Northern Adelbert, and Wikipedia treats its Trans-New Guinea affiliation as uncertain.

\section{Primary Educational Resources}
No primers, textbooks or school programmes for Usan were found in the search results.

\section{Linguistic and Reference Materials}
\begin{itemize}
\item Wikipedia, Usan language: \url{https://en.wikipedia.org/wiki/Usan_language}
\item Hein Reesink, \emph{Structures and their Functions in Usan} (1987, Studies in Language Companion Series 13), the standard grammar. The search results show only its publisher page: \url{https://www.degruyterbrill.com/document/doi/10.1075/slcs.13/html}
\item ASJP wordlist, Usan (a short comparative word list only): \url{https://asjp.clld.org/languages/USAN}
\item ASJP wordlist, Wanuma: \url{https://asjp.clld.org/languages/WANUMA}
\end{itemize}
Pick's 2020 dissertation on Proto-Northern Adelbert phonology was mentioned in a search summary, but I saw no URL for it, so it is not in the json list.

\section{Religious and Community Resources}
\begin{itemize}
\item Joshua Project, Usan, Wanuma in Papua New Guinea: \url{https://joshuaproject.net/people_groups/15743}
\end{itemize}

\section{Digital and Audio Resources}
No audio, video or corpus resources for Usan were found in the search results.

\section{Curriculum Development Context}
Usan is spoken by all adults but not by all young people, and it is not known to be taught in schools. Reesink's grammar is a solid linguistic base for orthography and lesson design. Reading materials, recordings and community input would all have to be created.
\end{document}